% ID: SEL-D0050-OA-E02
% ESTADO: validada. S051.
\textbf{a)} La derivada es
\[
B'(t)=\tfrac34(t-2)(t-6).
\]
Es positiva en $(0,2)$ y $(6,8)$ y negativa en $(2,6)$. Así, $2$ es un máximo relativo interior y $6$ un mínimo relativo interior. Para los extremos absolutos incluimos los extremos del dominio:
\[
\begin{array}{c|rrrr}
t&0&2&6&8\\\hline
B(t)&0&8&0&8
\end{array}
\]
El máximo absoluto es $8$ millones en $t=2$ y $t=8$. El mínimo absoluto es $0$ en $t=0$ y $t=6$.

\textbf{b)} El beneficio sube desde cero hasta los $8$ millones en los dos primeros años, desciende a cero entre el segundo y el sexto, y vuelve a subir hasta $8$ millones al finalizar el octavo año. La gráfica muestra también el cambio de concavidad en $t=4$, donde $B''(t)=\tfrac32t-6$ cambia de signo.
\begin{center}
\setlength{\unitlength}{0.9mm}
\begin{picture}(103,54)(-9,-7)
\qbezier(0,0)(45,0)(90,0)
\qbezier(0,0)(0,23)(0,46)
\put(92,0){$t$}
\put(2,44){$B(t)$}
\qbezier(0,0)(2.5,11.2109)(5,18.9062)
\qbezier(5,18.9062)(7.5,26.6016)(10,31.25)
\qbezier(10,31.25)(12.5,35.8984)(15,37.9688)
\qbezier(15,37.9688)(17.5,40.0391)(20,40)
\qbezier(20,40)(22.5,39.9609)(25,38.2812)
\qbezier(25,38.2812)(27.5,36.6016)(30,33.75)
\qbezier(30,33.75)(32.5,30.8984)(35,27.3438)
\qbezier(35,27.3438)(37.5,23.7891)(40,20)
\qbezier(40,20)(42.5,16.2109)(45,12.6562)
\qbezier(45,12.6562)(47.5,9.10156)(50,6.25)
\qbezier(50,6.25)(52.5,3.39844)(55,1.71875)
\qbezier(55,1.71875)(57.5,0.0390625)(60,0)
\qbezier(60,0)(62.5,-0.0390625)(65,2.03125)
\qbezier(65,2.03125)(67.5,4.10156)(70,8.75)
\qbezier(70,8.75)(72.5,13.3984)(75,21.0938)
\qbezier(75,21.0938)(77.5,28.7891)(80,40)
\put(-6,39){\scriptsize $8$}
\put(19,-5){\scriptsize $2$}
\put(39,-5){\scriptsize $4$}
\put(59,-5){\scriptsize $6$}
\put(79,-5){\scriptsize $8$}
\end{picture}
\end{center}
